\inicializarSeccion{Introducción}

\tab El juego mecánico de la torre de caída son atracciones diseñadas para elevar una góndola a cierta altura y después dejarla caer; lo que genera una sensación de caída libre seguida de una desaceleración controlada. En el reto analizamos el frenado magnético que es generado por las corrientes de Eddy, siendo una alternativa a utilizar el frenado por fricción entre superficies, lo cual disminuye en gran medida el desgaste mecánico y genera un frenado más suave.  Este tipo de frenado está basado en la inducción electromagnética; es cuando un conductor se mueve dentro del campo magnético, lo cual genera corrientes inducidas, las cuales producen una fuerza opuesta al movimiento \cite{libretexts2025eddy}.

\begin{figure}[H]
    \centering
    \includegraphics[width=0.5\textwidth]{\nat{parte1/ejemploFrenadoMagnetico}}
    \caption{Esquema de una torre de caída con frenado magnético.}
    \label{fig:ejemploFrenadoMagnetico}
\end{figure}

\tab Este proyecto se centra en calcular y simular la desaceleración de una góndola de una torre de caída libre utilizando el software MATLAB. Para lograr el objetivo de simular la caída libre, tenemos que tomar en cuenta dos etapas: en la primera, la caída libre, y en la segunda etapa, donde entra el frenado magnético, modelado como una fuerza proporcional a la velocidad. A partir del modelo planteado, se obtendrán las gráficas de velocidad y de aceleración en función del tiempo y de la altura, esto con el objetivo de analizar si el sistema puede frenar la góndola de una manera más segura para los usuarios del juego.

\tab En general, el objetivo de este proyecto nos permitirá entender cómo es que las leyes de Newton, Faraday y Lenz se relacionan entre sí y en un sistema ingenieril real \cite{physlabs2024magnetic}. Las simulaciones nos ayudan a evaluar los parámetros, como lo son la altura total de la torre, las zonas conductoras del frenado, la velocidad final y los posibles valores peligrosos de aceleración para los usuarios del juego. Este análisis es importante, ya que los frenos magnéticos se usan en sistemas de transporte y de otras atracciones, esto debido a que pueden funcionar sin fricción; hasta en algunos diseños no requieren una fuente externa de energía \cite{forrister2019how}. 
